\chapter{Middle convolution: dependence on parameters}
\label{chap:4}
\setcounter{section}{-1}

\section{Good schemes}\label{sec:4.0}

\sourcepara{4.0.1}
Recall that a ``good scheme'' is one which admits a map \(f\) of finite type to
a scheme \(T\) which is regular of dimension at most one. For variable good
schemes \(X\), and \(\ell\) any fixed prime number, we can speak of the
triangulated categories
\(\mathrm{D}^b_c(X[1/\ell],\overline{\mathbb{Q}}_\ell)\), which admit the full
Grothendieck formalism of the ``six operations''
(cf. \cite{De-ThFin}, \cite{De-WeilII}, \cite{Ek}, \cite{Me-SO}).

\section{The basic setting}\label{sec:4.1}

\sourcepara{4.1.1}
Throughout this section, we fix a prime number \(\ell\), and work over a
ground-ring \(R\) which is a normal noetherian integral domain in which \(\ell\)
is invertible, and such that \(\operatorname{Spec}(R)\) is a good scheme. We
work on
\(\mathbb{A}^1_R:=\operatorname{Spec}(R[x])\). We fix a monic polynomial
\(D(x)\) in \(R[x]\) of some degree \(d\geq1\) whose discriminant \(\Delta\) is
a unit in \(R\), and denote by \(D\subset\mathbb{A}^1_R\) the divisor defined
by the vanishing of \(D(x)\). We further assume that the polynomial \(D(x)\)
factors completely in \(R[x]\), say
\[
  D(x)=\prod_{1\leq i\leq d}(x-a_i)
  \quad\text{with }a_i-a_j\in R^\times\text{ for all }i\ne j.
\]
(This ``further assumption'' always holds after replacing
\(\operatorname{Spec}(R)\) by a connected, finite étale covering of itself.)
Then \(D=\coprod_i D_i\) is the disjoint union of the sections \(D_i\) defined
by the vanishing of \(x-a_i\).

\sourcepara{4.1.2}
We say that an object \(K\) in
\(\mathrm{D}^b_c(\mathbb{A}^1_R,\overline{\mathbb{Q}}_\ell)\) is adapted to
the stratification
\((\mathbb{A}^1_R-D,D=\coprod_iD_i)\) of \(\mathbb{A}^1_R\), if each of its
cohomology sheaves is lisse when restricted either to \(\mathbb{A}^1_R-D\) or
to any \(D_i\).

\sourcepara{4.1.3}
We say that an object \(K\) in
\(\mathrm{D}^b_c(\mathbb{A}^1_R,\overline{\mathbb{Q}}_\ell)\) is fibrewise
perverse if its restriction to each geometric fibre of \(\mathbb{A}^1_R\) over
\(\operatorname{Spec}(R)\) is perverse, i.e. if for any algebraically closed
field \(k\), and any ring homomorphism \(\varphi:R\to k\), the inverse image
\(K_\varphi\) of \(K\) in
\(\mathrm{D}^b_c(\mathbb{A}^1_k,\overline{\mathbb{Q}}_\ell)\) is perverse on
\(\mathbb{A}^1_k\).

\sourcepara{4.1.4}
We say that an object \(K\) in
\(\mathrm{D}^b_c(\mathbb{A}^1_R,\overline{\mathbb{Q}}_\ell)\) is fibrewise
tame if for any algebraically closed field \(k\), and any ring homomorphism
\(\varphi:R\to k\), the inverse image \(K_\varphi\) of \(K\) in
\(\mathrm{D}^b_c(\mathbb{A}^1_k,\overline{\mathbb{Q}}_\ell)\) is tame on
\(\mathbb{A}^1_k\) in the sense that for any dense open set
\(U\subset\mathbb{A}^1_k\) on which \(K_\varphi\) is lisse, each of the
cohomology sheaves \(\mathcal{H}^i(K_\varphi)\vert U\) is tamely ramified at
each point of \(\mathbb{P}^1_k-U\).

\section{Basic results in the basic setting}\label{sec:4.2}

\sourceheading[prop:4.2.1]{Proposition 4.2.1}
Hypotheses and notations as in \hyperref[par:4.1.1]{(4.1.1)} above, let \(K\)
in \(\mathrm{D}^b_c(\mathbb{A}^1_R,\overline{\mathbb{Q}}_\ell)\) be adapted to
\((\mathbb{A}^1_R-D,D=\coprod_iD_i)\). If the fraction field of \(R\) has
characteristic zero, then \(K\) is fibrewise tame.

\begin{proof}
\cite[4.7.1 and Remarque (ii)]{Ka-SE}, \cite[SGA 1, Exposé XIII, 5.5]{SGA}.
\end{proof}

\sourceheading[prop:4.2.2]{Proposition 4.2.2}
Hypotheses and notations as in \hyperref[par:4.1.1]{(4.1.1)} above, suppose
\(K\) in \(\mathrm{D}^b_c(\mathbb{A}^1_R,\overline{\mathbb{Q}}_\ell)\) is
adapted to \((\mathbb{A}^1_R-D,D=\coprod_iD_i)\) and is fibrewise tame. Then
the following conditions are equivalent:
\begin{enumerate}[label=(\arabic*)]
\item \(K\) is fibrewise perverse.
\item There exists a geometric fibre on which \(K\) is perverse.
\item Each of the following three conditions holds.
  \begin{enumerate}[label=(\roman*)]
  \item \(\mathcal{H}^i(K)\) vanishes for \(i\) outside \(\{0,-1\}\),
  \item \(\mathcal{H}^0(K)\vert\mathbb{A}^1_R-D\) vanishes, and
  \item denoting by
  \(j:\mathbb{A}^1_R-D\to\mathbb{A}^1_R\) the inclusion, the natural
  adjunction map
  \[
    \mathcal{H}^{-1}(K)\to j_*j^*\mathcal{H}^{-1}(K)
  \]
  is injective.
  \end{enumerate}
\end{enumerate}

\begin{proof}
Fix an algebraically closed field \(k\), and a ring homomorphism
\(\varphi:R\to k\). The inverse image \(K_\varphi\) of \(K\) on
\(\mathbb{A}^1_k\) is perverse if the conditions of \((3)\) above hold after the
base change \(\varphi:R\to k\); let us call these conditions \(3(\varphi)\).

For each integer \(i\), each of the sheaves
\(\mathcal{H}^i(K)\vert\mathbb{A}^1_R-D\) and
\(\mathcal{H}^i(K)\vert D\) is lisse. So for any \(\varphi:R\to k\) as above,
\((3\mathrm{i})\) \(\Leftrightarrow 3(\varphi\mathrm{i})\), and
\((3\mathrm{ii})\) \(\Leftrightarrow 3(\varphi\mathrm{ii})\). It is in showing
\((3\mathrm{iii})\) \(\Leftrightarrow 3(\varphi\mathrm{iii})\) that we make essential use
of the fibrewise tameness.

The key point \cite[4.7.2 and 4.7.3]{Ka-SE} is that for \(\mathcal{F}\) a
\(\overline{\mathbb{Q}}_\ell\)-sheaf on \(\mathbb{A}^1_R\) which is adapted to
\((\mathbb{A}^1_R-D,D=\coprod_iD_i)\) and fibrewise tame, we have:
\begin{enumerate}[label=(\arabic*)]
\item formation of \(j_*j^*\mathcal{F}\) commutes with arbitrary change of base
on \(\operatorname{Spec}(R)\), and
\item the sheaf
\[
  \mathcal{F}_{\mathrm{pct}}:=\operatorname{Ker}(\mathcal{F}\to j_*j^*\mathcal{F})
\]
is concentrated on \(D\), lisse on \(D\) (i.e., lisse on each \(D_i\)), and of
formation compatible with arbitrary change of base on \(\operatorname{Spec}(R)\).
\end{enumerate}
Applying this to \(\mathcal{H}^{-1}(K)\), we see that
\((3\mathrm{iii})\) \(\Leftrightarrow 3(\varphi\mathrm{iii})\).

Fixing a single \(\varphi\), we get \((2)\) \(\Rightarrow\) \((3)\). Varying
\(\varphi\), we get \((3)\) \(\Rightarrow\) \((1)\). And \((1)\) \(\Rightarrow\) \((2)\) is
trivial.
\end{proof}

\sourceheading[prop:4.2.3]{Proposition 4.2.3}
Hypotheses and notations as in \hyperref[par:4.1.1]{(4.1.1)} above, let \(K\)
in \(\mathrm{D}^b_c(\mathbb{A}^1_R,\overline{\mathbb{Q}}_\ell)\) be adapted to
\((\mathbb{A}^1_R-D,D=\coprod_iD_i)\), fibrewise perverse and fibrewise tame.
Then the following conditions are equivalent:
\begin{enumerate}[label=(\arabic*)]
\item on each geometric fibre, \(K\) is the middle extension of its restriction
to \(\mathbb{A}^1_k-D_k\).
\item there exists a geometric fibre on which \(K\) is the middle extension of
its restriction to \(\mathbb{A}^1_k-D_k\).
\item \(\mathcal{H}^0(K)=0\), and denoting by
\(j:\mathbb{A}^1_R-D\to\mathbb{A}^1_R\) the inclusion, the adjunction map is an
isomorphism
\[
  \mathcal{H}^{-1}(K)\cong j_*j^*\mathcal{H}^{-1}(K).
\]
\end{enumerate}

\begin{proof}
The kernel and cokernel of the adjunction map
\[
  \mathcal{H}^{-1}(K)\to j_*j^*\mathcal{H}^{-1}(K)
\]
are concentrated on \(D\), lisse on \(D\), and of formation compatible with
arbitrary change of base on \(\operatorname{Spec}(R)\), cf.
\cite[4.7.2--3]{Ka-SE}. Fixing a single \(\varphi\), we get
\((2)\) \(\Rightarrow\) \((3)\). Varying \(\varphi\), we get \((3)\) \(\Rightarrow\) \((1)\).
And \((1)\) \(\Rightarrow\) \((2)\) is trivial.
\end{proof}

\sourceheading[prop:4.2.4]{Proposition 4.2.4}
Hypotheses and notations as in \hyperref[par:4.1.1]{(4.1.1)} above, let \(K\)
in \(\mathrm{D}^b_c(\mathbb{A}^1_R,\overline{\mathbb{Q}}_\ell)\) be adapted to
\((\mathbb{A}^1_R-D,D=\coprod_iD_i)\), fibrewise perverse and fibrewise tame.
Denote by \(n\) the rank of the lisse sheaf
\(\mathcal{H}^{-1}(K)\vert\mathbb{A}^1_R-D\). Fix a subgroup
\(\Gamma\) of \(\operatorname{GL}(n,\overline{\mathbb{Q}}_\ell)\). Then the
following conditions are equivalent:
\begin{enumerate}[label=(\arabic*)]
\item on each geometric fibre, \(K\) is the middle extension of its restriction
to \(\mathbb{A}^1_k-D_k\), and the image of
\(\pi_1(\mathbb{A}^1_k-D_k,\text{base point})\) in
\(\operatorname{GL}(n,\overline{\mathbb{Q}}_\ell)\) under the monodromy
representation of
\(\mathcal{H}^{-1}(K)\vert\mathbb{A}^1_k-D_k\) is conjugate to \(\Gamma\).
\item On some geometric fibre, \(K\) is the middle extension of its restriction
to \(\mathbb{A}^1_k-D_k\), and the image of
\(\pi_1(\mathbb{A}^1_k-D_k,\text{base point})\) in
\(\operatorname{GL}(n,\overline{\mathbb{Q}}_\ell)\) under the monodromy
representation of
\(\mathcal{H}^{-1}(K)\vert\mathbb{A}^1_k-D_k\) is conjugate to \(\Gamma\).
\end{enumerate}

\begin{proof}
In view of the preceding \hyperref[prop:4.2.3]{Proposition~4.2.3}, this results from the fact that for a lisse
\(\overline{\mathbb{Q}}_\ell\)-sheaf \(\mathcal{F}\) on
\(\mathbb{A}^1_R-D\) (namely
\(\mathcal{H}^{-1}(K)\vert\mathbb{A}^1_R-D\)) of rank denoted \(n\) which is
fibrewise tame, the function on geometric points of \(\operatorname{Spec}(R)\)
given by
\[
  (\varphi:R\to k)\longmapsto
  \text{the conjugacy class in }\operatorname{GL}(n,\overline{\mathbb{Q}}_\ell)
  \text{ of the image of }
  \pi_1(\mathbb{A}^1_k-D_k,\text{any base point})
\]
is constant (cf. \cite[8.17.13]{Ka-ESDE}).
\end{proof}

\sourceheading[cor:4.2.5]{Corollary 4.2.5}
Hypotheses and notations as in \hyperref[par:4.1.1]{(4.1.1)} above, let \(K\)
in \(\mathrm{D}^b_c(\mathbb{A}^1_R,\overline{\mathbb{Q}}_\ell)\) be adapted to
\((\mathbb{A}^1_R-D,D=\coprod_iD_i)\), fibrewise perverse and fibrewise tame.
Then the following conditions are equivalent:
\begin{enumerate}[label=(\arabic*)]
\item on each geometric fibre, \(K\cong\overline{\mathbb{Q}}_\ell[1]\).
\item there exists a geometric fibre on which
\(K\cong\overline{\mathbb{Q}}_\ell[1]\).
\end{enumerate}

\begin{proof}
This is the case \(n=1\), \(\Gamma=\{1\}\) of
\hyperref[prop:4.2.4]{Proposition~4.2.4}.
\end{proof}

\sourceheading[cor:4.2.6]{Corollary 4.2.6}
Hypotheses and notations as in \hyperref[par:4.1.1]{(4.1.1)} above, let \(K\)
in \(\mathrm{D}^b_c(\mathbb{A}^1_R,\overline{\mathbb{Q}}_\ell)\) be adapted to
\((\mathbb{A}^1_R-D,D=\coprod_iD_i)\), fibrewise perverse and fibrewise tame.
Then the following conditions are equivalent:
\begin{enumerate}[label=(\arabic*)]
\item on each geometric fibre, \(K\) is perverse irreducible and is the middle
extension of its restriction to \(\mathbb{A}^1_k-D_k\).
\item there exists a geometric fibre on which \(K\) is perverse irreducible and
is the middle extension of its restriction to \(\mathbb{A}^1_k-D_k\).
\end{enumerate}

\begin{proof}
This is the case ``\(\Gamma\) an irreducible subgroup of
\(\operatorname{GL}(n,\overline{\mathbb{Q}}_\ell)\)'' of
\hyperref[prop:4.2.4]{Proposition~4.2.4}, for all such \(\Gamma\).
\end{proof}

\sourceheading[prop:4.2.7]{Proposition 4.2.7}
Hypotheses and notations as in \hyperref[par:4.1.1]{(4.1.1)} above, let \(K\)
in \(\mathrm{D}^b_c(\mathbb{A}^1_R,\overline{\mathbb{Q}}_\ell)\) be adapted to
\((\mathbb{A}^1_R-D,D=\coprod_iD_i)\), fibrewise perverse and fibrewise tame.
Denote by \(j:\mathbb{A}^1_R-D\to\mathbb{A}^1_R\) the inclusion. Then the
following conditions are equivalent:
\begin{enumerate}[label=(\arabic*)]
\item on each geometric fibre,
\(K\cong j_*\mathcal{L}_{\chi}(x-\alpha)[1]\) for some nontrivial Kummer sheaf
\(\mathcal{L}_\chi\) and some \(\alpha\).
\item there exists a geometric fibre on which
\(K\cong j_*\mathcal{L}_{\chi}(x-\alpha)[1]\) for some nontrivial Kummer sheaf
and some \(\alpha\).
\item Each of the following conditions holds:
  \begin{enumerate}[label=(\roman*)]
  \item \(\mathcal{H}^0(K)=0\),
  \item \(\mathcal{H}^{-1}(K)\vert\mathbb{A}^1_R-D\) is lisse of rank one,
  \item \(\mathcal{H}^{-1}(K)\cong j_*j^*\mathcal{H}^{-1}(K)\),
  \item \(\mathcal{H}^{-1}(K)\vert D_i\) is lisse of rank one for all but
  exactly one value \(i_0\) of \(i\), and for this value
  \(\mathcal{H}^{-1}(K)\vert D_{i_0}=0\).
  \end{enumerate}
\end{enumerate}

\begin{proof}
It is clear that \((1)\) \(\Rightarrow\) \((2)\) and \((2)\) \(\Rightarrow\) \((3)\). To see
that \((3)\) \(\Rightarrow\) \((1)\), we may reduce by an additive translation on
\(\mathbb{A}^1_R\) to the case where the divisor \(D_{i_0}\) is the divisor
\(x=0\). In this case, we must show that for any geometric point
\(\varphi:R\to k\), the only lisse rank one
\(\overline{\mathbb{Q}}_\ell\)-sheaf on \(\mathbb{G}_{m,k}\) which is ramified
at \(0\) and which is tame at both \(0\) and \(\infty\) is a nontrivial Kummer
sheaf \(\mathcal{L}_\chi\). But this is a tautology.
\end{proof}

\sourceheading[prop:4.2.8]{Proposition 4.2.8}
Hypotheses and notations as in \hyperref[par:4.1.1]{(4.1.1)} above, let \(K\)
in \(\mathrm{D}^b_c(\mathbb{A}^1_R,\overline{\mathbb{Q}}_\ell)\) be adapted to
\((\mathbb{A}^1_R-D,D=\coprod_iD_i)\), fibrewise perverse and fibrewise tame.
Then the following conditions are equivalent:
\begin{enumerate}[label=(\arabic*)]
\item on each geometric fibre, \(K\cong\delta_\alpha\) for some \(\alpha\).
\item there exists a geometric fibre on which \(K\cong\delta_\alpha\) for some
\(\alpha\).
\item \(\mathcal{H}^{-1}(K)=0\), and
\(\mathcal{H}^0(K)\vert D_i=0\) for all but exactly one value \(i_0\) of \(i\),
and for this value \(\mathcal{H}^{-1}(K)\vert D_{i_0}\) is lisse of rank one.
\end{enumerate}

\begin{proof}
It is clear that \((1)\) \(\Rightarrow\) \((2)\) and \((2)\) \(\Rightarrow\) \((3)\), and
\((3)\) \(\Rightarrow\) \((1)\) is obvious as well.
\end{proof}

\section{Middle convolution in the basic setting}\label{sec:4.3}

\sourcepara{4.3.1}
We continue to work in the basic setting \hyperref[par:4.1.1]{(4.1.1)}. Fix a
nontrivial character \(\chi\) of the group
\[
  \pi_1^{\mathrm{tame}}(\mathbb{G}_{m,R})
  =
  \prod_{\ell\text{ inv in }R}\mathbb{Z}_\ell(1).
\]

\sourcepara{4.3.2}
Denote by \(\mathcal{C}(R,D)\) the full subcategory of
\(\mathrm{D}^b_c(\mathbb{A}^1_R,\overline{\mathbb{Q}}_\ell)\) consisting of
all objects \(K\) which are adapted to
\((\mathbb{A}^1_R-D,D=\coprod_iD_i)\), fibrewise perverse and fibrewise tame.
We will now define a middle convolution functor from \(\mathcal{C}(R,D)\) to
itself,
\[
  K\longmapsto K*_{\mathrm{mid}+}j_*\mathcal{L}_\chi[1],
\]
whose formation commutes with arbitrary change of base on \(R\), and which for
\(R\) a field of characteristic \(\ne\ell\) coincides with its namesake
\hyperref[par:2.6.2]{(2.6.2)}. It is defined via the relative compactification
of the map \(\operatorname{pr}_2\) as in
\hyperref[par:2.8.3]{(2.8.3)}--\hyperref[prop:2.8.4]{Proposition~2.8.4},
\[
  j:\mathbb{A}^1_x\times\mathbb{A}^1_t
  \hookrightarrow
  \mathbb{P}^1_x\times\mathbb{A}^1_t,
\]
by extending \(K_x\otimes(j_*\mathcal{L}_\chi[1])_{t-x}\) across
\(Z:=\infty_{\mathbb{A}^1}\) by
\(\tau^Z_{\leq -2}\mathrm{R}j_*\), \cite[1.4.13]{BBD} and then taking
\(\mathrm{R}\operatorname{pr}_{2*}\) as in the earlier discussion. The key
point is that
\(\mathrm{R}j_*(K_x\otimes(j_*\mathcal{L}_\chi[1])_{t-x})\vert Z\) is lisse,
and of formation compatible with change of base on \(R\) (by
\hyperref[lem:4.3.8]{Lemma~4.3.8} below). Therefore the object
\[
  \tau^Z_{\leq -2}
  \mathrm{R}j_*(K_x\otimes(j_*\mathcal{L}_\chi[1])_{t-x})\vert Z
\]
is lisse, and of formation compatible with change of base on \(R\).

\sourcepara{4.3.3}
Since the formation of \(\mathrm{R}\operatorname{pr}_{2*}\) commutes with
passage to fibres, by proper base change, we know that
\(K*_{\mathrm{mid}+}j_*\mathcal{L}_\chi[1]\) is fibrewise perverse, of
formation compatible with arbitrary change of base on \(R\), and fibre by
fibre equal to its namesake.

\sourcepara{4.3.4}
We also have an explicit triangle relating this middle convolution to the
\(!\) convolution and to a lisse sheaf on \(\mathbb{A}^1_R\). On
\(\mathbb{P}^1_R\times\mathbb{A}^1_R\) we have a distinguished triangle
\[
\begin{aligned}
  &\tau^Z_{\leq -3}
  \mathrm{R}j_*(K_x\otimes(j_*\mathcal{L}_\chi[1])_{t-x}) \\
  &\qquad\to
  \tau^Z_{\leq -2}
  \mathrm{R}j_*(K_x\otimes(j_*\mathcal{L}_\chi[1])_{t-x}) \\
  &\qquad\to
  j_*\bigl(\mathcal{H}^{-2}(K_x\otimes(j_*\mathcal{L}_\chi[1])_{t-x})[2]\bigr)
  \vert Z.
\end{aligned}
\]

\sourcepara{4.3.5}
But
\(\tau^Z_{\leq -3}\mathrm{R}j_*(K_x\otimes(j_*\mathcal{L}_\chi[1])_{t-x})\) is
\(j_!(K_x\otimes(j_*\mathcal{L}_\chi[1])_{t-x})\), so rotating this triangle
gives a distinguished triangle
\[
\begin{aligned}
  &j_*\bigl(\mathcal{H}^{-2}(K_x\otimes(j_*\mathcal{L}_\chi[1])_{t-x})[1]\bigr)
  \vert Z \\
  &\qquad\to
  j_!(K_x\otimes(j_*\mathcal{L}_\chi[1])_{t-x}) \\
  &\qquad\to
  \tau^Z_{\leq -2}\mathrm{R}j_*(K_x\otimes(j_*\mathcal{L}_\chi[1])_{t-x}),
\end{aligned}
\]
in which the first term is (a lisse sheaf on \(Z\))[1].

\sourcepara{4.3.6}
Applying \(\mathrm{R}\operatorname{pr}_{2*}\) gives a distinguished triangle on
\(\mathbb{A}^1_R\)
\[
  (\text{a lisse sheaf on }\mathbb{A}^1_R)[1]
  \to
  K*_{!+}j_*\mathcal{L}_\chi[1]
  \to
  K*_{\mathrm{mid}+}j_*\mathcal{L}_\chi[1].
\]
Using this triangle, we see the adaptedness of
\(K*_{\mathrm{mid}+}j_*\mathcal{L}_\chi[1]\) to the same stratification
\((\mathbb{A}^1_R-D,D=\coprod_iD_i)\) to which \(K\) was adapted.

\sourcepara{4.3.7}
We know fibre by fibre that tameness is preserved: this is automatic for
characteristic zero fibres, and for characteristic \(p>0\) fibres it was proven
in \hyperref[cor:3.3.6]{Corollary~3.3.6}.

\sourceheading[lem:4.3.8]{Lemma 4.3.8}
Let \(S\) be an irreducible noetherian scheme, \(X/S\) smooth, and \(D\) in
\(X\) a smooth/\(S\) divisor. For \(\mathcal{F}\) lisse on \(X-D\) and tame
along \(D\), \(j:X-D\to X\) and \(i:D\to X\) the inclusions, we have
\begin{enumerate}[label=(\arabic*)]
\item formation of \(j_*\mathcal{F}\) and of
\(\mathrm{R}j_*\mathcal{F}\) on \(X\) commutes with arbitrary change of base on
\(S\),
\item the sheaf \(i^*j_*\mathcal{F}\) on \(D\) is lisse, and formation of
\(i^*j_*\mathcal{F}\) on \(D\) commutes with arbitrary change of base on \(S\).
\end{enumerate}

\begin{proof}
\cite[4.7.2 and 4.7.3]{Ka-SE}.
\end{proof}

\sourceheading[prop:4.3.9]{Proposition 4.3.9}
Hypotheses and notations as in \hyperref[par:4.1.1]{(4.1.1)} above, let \(K\)
in \(\mathrm{D}^b_c(\mathbb{A}^1_R,\overline{\mathbb{Q}}_\ell)\) be adapted to
\((\mathbb{A}^1_R-D,D=\coprod_iD_i)\), fibrewise perverse and fibrewise tame.
If \(K\) is fibrewise a middle extension, then its fibrewise index of rigidity
is constant.

\begin{proof}
Consider the lisse tame sheaf on \(\mathbb{P}^1_R-\{D\coprod\infty\}\),
\[
  \mathcal{F}:=\operatorname{End}\bigl(
  \mathcal{H}^{-1}\vert\mathbb{P}^1_R-\{D\coprod\infty\}\bigr).
\]
Its \(j_*\mathcal{F}\) on \(\mathbb{P}^1_R\) is adapted to
\((\mathbb{P}^1_R-\{D\coprod\infty\},\{D\coprod\infty\})\), fibrewise tame, and
of formation compatible with arbitrary change of base on \(\operatorname{Spec}(R)\),
by \hyperref[lem:4.3.8]{Lemma~4.3.8}. By the Euler--Poincaré formula, the fibrewise
\(\chi\) of \(j_*\mathcal{F}\) is constant. Indeed, if
\[
\begin{aligned}
  n&:=\operatorname{rank}\text{ of }\mathcal{F}
  \text{ on }\mathbb{P}^1_R-\{D\coprod\infty\},\\
  n_\infty&:=\operatorname{rank}\text{ of }j_*\mathcal{F}\vert\infty_R,\\
  n_i&:=\operatorname{rank}\text{ of }j_*\mathcal{F}\vert D_i,
  \text{ for each of the }d\text{ sections }D_i,
\end{aligned}
\]
then the fibrewise \(\chi\) of \(j_*\mathcal{F}\), i.e., the fibrewise index of
rigidity, is
\[
  (1-d)n+n_\infty+\sum_i n_i.
\]
\end{proof}

\sourceheading[thm:4.3.10]{Theorem 4.3.10}
Hypotheses and notations as in \hyperref[par:4.1.1]{(4.1.1)} above, let \(K\)
in \(\mathrm{D}^b_c(\mathbb{A}^1_R,\overline{\mathbb{Q}}_\ell)\) be adapted to
\((\mathbb{A}^1_R-D,D=\coprod_iD_i)\), fibrewise perverse and fibrewise tame.
Suppose \(K\) is fibrewise perverse irreducible of type \((2d)\) in the sense
of \hyperref[thm:3.3.3]{Theorem~3.3.3}. Then both \(K\) and
\(K*_{\mathrm{mid}+}j_*\mathcal{L}_\chi[1]\) are fibrewise perverse
irreducible of type \((2d)\) and have the same index of rigidity.

\begin{proof}
Suppose first that \(R\) is not a \(\mathbb{Q}\)-algebra. Then
\(\operatorname{Spec}(R)\) has points with finite residue characteristic \(p\).
Once we know that
\[
  K\longmapsto K*_{\mathrm{mid}+}j_*\mathcal{L}_\chi[1]
\]
is stable on fibrewise perverses which are fibrewise tame and adapted to \(D\),
we just look fibrewise (legitimate by the earlier fibrewise results) and take a
fibre in characteristic \(p>0\), where we can use
\hyperref[thm:3.3.3]{Theorem~3.3.3(3)}.

How do we get this result if we do not assume \(R\) has points in some positive
characteristic? By passing to fibres, we first reduce to the case when \(R\) is
itself an algebraically closed field of characteristic zero. Since \(\pi_1\) of
open curves in characteristic zero does not ``see'' extensions of algebraically
closed fields, we may assume that our field \(R\) is ``just'' an algebraic
closure of the finitely generated field
\(\mathbb{Q}(\text{coef's of }D(x))\), and try working over the ring
\[
  R_0:=\mathbb{Z}[\text{coef's of }D(x),1/\Delta].
\]
This ring \(R_0\) is way too big, but it does have points of characteristic
\(p\) for all \(p\gg0\). Our
\(\mathcal{H}^{-1}\vert\mathbb{A}^1_R-D\) has
\(\operatorname{GL}(n,\mathcal{O}_\lambda)\) form, \(\mathcal{O}_\lambda\)
denoting the ring of integers in a finite extension \(E_\lambda\) of
\(\mathbb{Q}_\ell\), with finite residue field \(\mathbb{F}_\lambda\). The key
point is that for \(p\gg0\) the group
\(\operatorname{GL}(n,\mathcal{O}_\lambda)\) is prime to \(p\), because it is
an extension of the finite group
\(\operatorname{GL}(n,\mathbb{F}_\lambda)\) by the pro-\(\ell\) group
\(1+\lambda M_n(\mathcal{O}_\lambda)\). For any \(p\gg0\), we can, by applying
\hyperref[lem:5.9.3]{Lemma~5.9.3} and then extending the residue field, embed our
\(R_0\) inside a complete discrete valuation ring \(R_1\) with algebraically
closed residue field of characteristic \(p\).

When we do this, \(D\) stays a good \(D\). By the specialization theorem for
the prime to \(p\) fundamental group of the open curve \(\mathbb{A}^1_{R_1}-D\),
we see that the sheaf
\(\mathcal{H}^{-1}\vert\mathbb{A}^1_R-D\) extends to a lisse sheaf on
\(\mathbb{A}^1_{R_1}-D\), which will be irreducible if and only if our
original \(\mathcal{H}^{-1}\vert\mathbb{A}^1_R-D\) was geometrically
irreducible. Denoting by
\[
  j:\mathbb{A}^1_{R_1}-D\to\mathbb{A}^1_{R_1}
\]
the inclusion, \(j_*(\mathcal{H}^{-1}\vert\mathbb{A}^1_{R_1}-D)\) is adapted,
fibrewise perverse, fibrewise tame, and fibrewise a middle extension. By
\hyperref[cor:4.2.6]{Corollary~4.2.6}, \hyperref[prop:4.2.7]{Proposition~4.2.7},
and \hyperref[cor:4.2.5]{Corollary~4.2.5}, the fibrewise ``type \(2d\)
irreducible'' condition is still satisfied by our extended \(K\) over \(R_1\).
Now we apply the already known case of the result over the mixed characteristic
\(R_1\) to get the index of rigidity on the general fibre.
\end{proof}

\sourceheading[thm:4.3.11]{Theorem 4.3.11}
Hypotheses and notations as in \hyperref[par:4.1.1]{(4.1.1)} above, let \(K\)
in \(\mathrm{D}^b_c(\mathbb{A}^1_R,\overline{\mathbb{Q}}_\ell)\) be adapted to
\((\mathbb{A}^1_R-D,D=\coprod_iD_i)\), fibrewise perverse and fibrewise tame.
Suppose \(K\) is fibrewise perverse irreducible of type \((2d)\) in the sense
of \hyperref[thm:3.3.3]{Theorem~3.3.3}. For
\(j:\mathbb{A}^1_R-D\to\mathbb{A}^1_R\) the inclusion, write
\[
  j^*K=\mathcal{F}[1],
  \qquad
  j^*(K*_{\mathrm{mid}+}j_*\mathcal{L}_\chi[1])=\mathcal{G}[1],
\]
with \(\mathcal{F}\) and \(\mathcal{G}\) lisse sheaves on
\(\mathbb{A}^1_R-D\). Then on every geometric fibre, the local monodromies of
\(\mathcal{F}\) and of \(\mathcal{G}\) are related as in
\hyperref[cor:3.3.6]{Corollary~3.3.6} and
\hyperref[cor:3.3.7]{Corollary~3.3.7}.

\begin{proof}
Denote by \(L\) the set of those prime numbers which are invertible on \(R\).
Because \(\mathcal{F}\) is lisse on
\(\mathbb{A}^1_R-\coprod D_i=\mathbb{P}^1_R-\{\infty_R,\coprod D_i\}\) and
tame along each of the ``missing'' sections \(Z:=\infty_R\) or \(D_i\), its
local monodromy along each of these sections is a pair
\[
  \bigl(
  \text{a lisse }\mathbb{Q}_\ell\text{-sheaf }\mathcal{F}[Z]\text{ on }Z,
  \text{ an action of }\mathbb{Z}_L(1)\text{ on }\mathcal{F}[Z]
  \bigr)
\]
whose formation is compatible with arbitrary change of base on
\(\operatorname{Spec}(R)\) (cf. \cite[1.7.8]{De-WeilII}, also
\cite[4.7.2]{Ka-SE}). Similarly for \(\mathcal{G}\). So in order the verify
the asserted relations between the fibrewise local monodromies of
\(\mathcal{F}\) and \(\mathcal{G}\), it suffices to do so on a single
geometric fibre.

If \(\operatorname{Spec}(R)\) has a point of residue characteristic \(p\), we
are done by \hyperref[cor:3.3.6]{Corollary~3.3.6} and
\hyperref[cor:3.3.7]{Corollary~3.3.7}. If not, we reduce first to the case when
\(R\) is an algebraically closed field of characteristic zero, and then we
specialize both \(\mathcal{F}\) and \(\mathcal{L}_\chi\) into characteristic
\(p\gg0\) by the argument of \hyperref[thm:4.3.10]{Theorem~4.3.10}, which reduces us to the case
when \(R\) is a mixed characteristic discrete valuation ring.
\end{proof}
